\section{Conclusion}
\label{sec:conclusion}

We measured what unsupervised domain adaptation adds to a simulation-trained lens finder on
Euclid Q1, at the false-positive rate a search runs at. With a class-balanced target, each of
four published objectives raised the TPR at FPR${}=10^{-3}$ from $\dFmBaseline{}$ to between
$\dHeadMin{}$ and $\dHeadMax{}$, measured against a control that differs from one of them only
in its adaptation weight. Because the target batches are balanced, that signal includes a known
class proportion as well as real images, and our design cannot separate the two. In a second
pre-registered experiment, a first test toward survey prevalence, we diluted the target to one
lens per ten galaxies and drew its batches without the labels. DANN still raised TPR over the
control, by $\dlGDann{}$ $[\dlGDannLo{}, \dlGDannHi{}]$, and BNM by $\dlGBnm{}$
$[\dlGBnmLoThree{}, \dlGBnmHi{}]$, though for both the comparison with the balanced target
cannot rule out that much of the original gain was lost; for CORAL and FreeMatch we cannot detect a gain. We
could not establish a ranking of CORAL, DANN and FreeMatch: at their declared read-offs every
paired interval includes zero, and an \rwFiveKWord{}-checkpoint average favours FreeMatch with an
interval excluding zero under only one of its two pre-declared score combinations. Settling the
ranking needs more graded lenses rather than another read-off convention. Two conditions limit how far
these numbers carry. Our simulation-only baseline is weak, so the gains are measured from a low
starting point. And every target we used, the diluted one included, is built from a lens
catalogue that a new survey will not have, and is far denser in lenses than the sky. The next
experiment is a target with no catalogue in it at all: random galaxies drawn from the survey, at
its own prevalence, which removes the remaining use of the labels and the remaining gap in lens
density.
